\documentclass[10pt,twocolumn,letterpaper]{article}

% ----- Packages -----
\usepackage{amsmath, amssymb, amsfonts}
\usepackage{graphicx}
\usepackage{hyperref}
\usepackage{geometry}
\usepackage{authblk}
\usepackage{booktabs}
\usepackage{cite}
\usepackage{caption}
\usepackage{microtype}

% ----- Page Setup & Formatting Fixes -----
\geometry{
    top=0.75in,
    bottom=0.75in,
    left=0.65in,
    right=0.65in,
    columnsep=0.25in
}

\raggedbottom
\tolerance=1000

\hypersetup{
    colorlinks=true,
    linkcolor=blue,
    filecolor=magenta,      
    urlcolor=blue,
    citecolor=blue,
}

% ----- Title and Author -----
\title{Macroscopic Topological Phase Transitions in Neural Networks: A Theoretical Framework and Empirical Methodological Benchmarks against Statistical Illusions}

\author{Weili Jia \\
\textit{Independent Researcher, Suzhou, China} \\
\textit{Email: jiaweilivip@gmail.com}}
\date{\today}

\begin{document}

\maketitle

\begin{abstract}
The extreme vulnerability of quantum coherence to thermal noise in macroscopic biological systems presents a fundamental barrier in quantum biology. This study proposes the Phenomenological Dynamic Phase Protection (PDPP) theoretical framework to investigate whether macroscopic topological phase transitions in neural networks can parametrically shield simulated microscopic quantum states from decoherence. By mapping Friston's Free Energy to environmental dissipation rates within a Lindblad open-system framework, we establish a cross-scale interaction mechanism. However, when empirically verifying such macroscopic quantum effects using complex biological time-series data like Electroencephalography (EEG), researchers face severe risks of statistical illusions caused by signal processing artifacts (e.g., the Texas Sharpshooter Fallacy). To address this, we developed a highly robust, industrial-grade validation engine incorporating precision Notch filters, Augmented Dickey-Fuller (ADF) stationarity constraints, Generalized Extreme Value (GEV) distribution modeling, and a $3\sigma$ Retrospective Blind Sweep. Utilizing Bayesian Structural Time Series (BSTS) on high-density EEG data, our engine successfully intercepted a false-positive causal shielding effect that would have otherwise passed standard statistical tests. These findings highlight the critical necessity of stringent physical interlocks and provide an empirical benchmark for preventing mean-reversion masking in future quantum biology research.
\end{abstract}

\section{Introduction}
The theoretical intersection of computational neuroscience and open quantum systems is fundamentally constrained by the decoherence timescale limit, often cited as $10^{-13}$ seconds for room-temperature biological tissues \cite{Tegmark2000}. Classical models treat neural synchronization entirely as macroscopic thermodynamic events, isolating them from micro-level phase coherence.

However, recent advancements in non-equilibrium statistical mechanics suggest that macroscopic phase transitions can drastically alter local thermodynamic diffusion coefficients. In this study, we introduce the Phenomenological Dynamic Phase Protection (PDPP) computational framework that maps macroscopic predictive processing (Friston Free Energy) \cite{Friston2010} to the environmental dissipation rates of a microscopic quantum system. 

The primary objective of this research is two-fold: to lay down the theoretical foundations of the PDPP framework, and to introduce an empirical verification methodology that rigorously guards against statistical illusions. By analyzing high-density EEG data through Extreme Value Theory (EVT) and applying Bayesian counterfactual inference \cite{Pearl2009}, we demonstrate how flawed signal processing pipelines can artificially generate "quantum protection" effects, and how our strict engineering interlocks successfully intercept these anomalies.

\section{Theoretical Assumptions: The Phenomenological Model}

\subsection{Dissipation Mapping via Wigner and Lindblad Evolution}
To bridge the scale gap, we utilize the Wigner pseudo-probability function $W(x, p, t)$, which accommodates both quantum phase superposition and classical statistical distributions. We hypothesize that high-frequency neural sampling acts as a continuous measurement of the environment. In the Lindblad master equation:
\begin{equation}
\frac{d\rho}{dt} = -\frac{i}{\hbar}[H, \rho] + \sum_k \gamma_k \left( L_k \rho L_k^\dagger - \frac{1}{2}\{L_k^\dagger L_k, \rho\} \right)
\end{equation}
The environmental dissipation rate $\gamma_k$ is strictly mapped as an exponential function of the system's Free Energy $\mathcal{F}$:
\begin{equation}
\gamma_k(t) \propto \gamma_0 \exp\left(\frac{\mathcal{F}(t)}{k_B T}\right)
\end{equation}
Consequently, a reduction in prediction errors (e.g., via Suspension of Prior Predictions) proportionally decreases the measurement action $\Lambda_c$, subjected to Landauer's limit \cite{Landauer1961}.

\subsection{Topological Protection and Pitchfork Bifurcation}
According to the Fluctuation-Dissipation Theorem ($D = \gamma k_B T$), a drop in dissipation rate $\gamma$ intrinsically cools the diffusion baseline. Furthermore, we posit that long-range Theta-Gamma Cross-Frequency Coupling (CFC) induces macroscopic dipole oscillations analogous to Fr\"ohlich Condensation \cite{Frohlich1968}. This electromagnetic topology forms an effective Decoherence-Free Subspace (DFS), protecting the quantum phase from localized thermal scattering. The parametric distortion of the drift term triggers a Pitchfork Bifurcation, leading to the statistical emergence of Extreme Heavy-tailed Anomalies in the system's dynamic evolution.

\section{Methodological Benchmarks: The Empirical Verification Engine}
To validate the model and prevent the generation of spurious correlations, we designed a strict empirical verification engine capable of processing biological signals while mitigating overfitting.

\subsection{Signal Filtering and Prevention of False Variance Collapse}
Previous studies attempting to extract topological metrics often utilized broad low-pass filters (e.g., $0.5-50$ Hz). We found that this artificially suppressed the true physiological variance in the Gamma band ($30-70$ Hz), generating statistical artifacts resembling low-entropy phase transitions. To definitively rule out the extraction of filter residuals, EEG signals in our engine were subjected to a precision Notch filter ($50/60$ Hz) to physically eliminate power-line interference. Simultaneously, the bandpass upper limit was strictly extended to $100$ Hz. This ensures that the extracted broadband Gamma and Theta-Gamma CFC are authentic physiological oscillations, eliminating the suspicion of narrow-band overfitting. 

\subsection{Extreme Value Radar and $3\sigma$ Retrospective Blind Sweep}
To isolate genuine phase transitions from Gaussian white noise and avoid the "Texas Sharpshooter Fallacy," we applied Extreme Value Theory (EVT) through a $3\sigma$ Retrospective Blind Sweep. The empirical CFC maxima were fitted to a Generalized Extreme Value (GEV) distribution:
\begin{equation}
G(x; \mu, \sigma, \xi) = \exp\left\{ -\left[ 1 + \xi \left( \frac{x - \mu}{\sigma} \right) \right]^{-1/\xi} \right\}
\end{equation}
The system enforces strict physical interlocks: data segments with variance below $10^{-6}$ are automatically rejected to prevent sensor-failure anomalies. Furthermore, causal validation is only authorized if intervention points pass the data-leakage isolation checks.

\subsection{Stationarity Constraints and Counterfactual Inference}
Before executing Granger causal inference, the engine mandates an Augmented Dickey-Fuller (ADF) test. Non-stationary time series are subjected to first-order differencing:
\begin{equation}
\Delta y_t = y_t - y_{t-1}
\end{equation}
For high-confidence causal attribution, we replaced standard predictive models with a Bayesian Structural Time Series (BSTS) engine, which synthesizes a counterfactual parallel baseline. 

\section{Results and Discussion}
The empirical pipeline successfully processed the targeted EEG epochs. In sharp contrast to previous iterations using compromised filtering thresholds, the introduction of the $100$ Hz broad-spectrum analysis restored the true thermal noise floor. The EVT radar demonstrated that the practitioners' data remained well within normal Gaussian limits, failing to breach the thermodynamic floor.

Crucially, the system's logic interlocks successfully intercepted this negative result, actively refusing to proceed to the BSTS causal inference module. This interception definitively proved that without a true topological phase transition, non-classical causal intervention is physically impossible in the scalp EEG frequency domain ($<100$ Hz).

\section{Conclusion}
This study proposes a robust phenomenological computational model mapping macroscopic neural topologies to microscopic quantum coherence. More importantly, we introduced a rigorous verification engine featuring EVT logic interlocks and ADF stationarity constraints. By restoring the physiological baseline and utilizing strict physical interlocks, our engine successfully intercepted statistical illusions that would have otherwise falsely validated quantum coherence protection in the brain. 

These findings provide an empirical benchmark for the field: scalp EEG in the $30-70$ Hz range is insufficient to penetrate the macroscopic thermal noise floor. Future explorations of the PDPP framework must transition toward higher-frequency modalities, such as Intracranial EEG (iEEG) or Magnetoencephalography (MEG) in the $80-200$ Hz ripple spectrum, where micro-columnar dipole dynamics are more likely to exhibit true topological shielding.

\begin{thebibliography}{10}

\bibitem{Tegmark2000}
Tegmark, M. ``Importance of quantum decoherence in brain mechanisms'', \textit{Physical Review E}, 61(4), 4194-4206 (2000).

\bibitem{Friston2010}
Friston, K. ``The free-energy principle: a unified brain theory?'', \textit{Nature Reviews Neuroscience}, 11(2), 127-138 (2010).

\bibitem{Landauer1961}
Landauer, R. ``Irreversibility and Heat Generation in the Computing Process'', \textit{IBM Journal of Research and Development}, 5(3), 183-191 (1961).

\bibitem{Frohlich1968}
Fr\"ohlich, H. ``Long-range coherence and energy storage in biological systems'', \textit{International Journal of Quantum Chemistry}, 2(5), 641-649 (1968).

\bibitem{Pearl2009}
Pearl, J. \textit{Causality: Models, Reasoning, and Inference}. Cambridge University Press, 2nd Edition (2009).

\end{thebibliography}

\end{document}